% Comandos de composição do deck: um pequeno sistema visual reutilizável.
\newcommand{\NstarHistorical}{\ensuremath{N^{\star}}}
\newcommand{\rankrho}{\ensuremath{\rho_{\mathrm{rank}}}}
\newcommand{\WA}{\textit{Winner Agreement}}
\newcommand{\Wmat}{\ensuremath{\mathbf{W}}}
\newcommand{\Rmat}{\ensuremath{\mathbf{R}}}
\newcommand{\AW}{\ensuremath{A_W}}
\newcommand{\AR}{\ensuremath{A_R}}
\newcommand{\SR}{\ensuremath{S_R}}
\newcommand{\DR}{\ensuremath{D_R}}
\newcommand{\PredictiveCooperationProfile}{perfil de cooperação preditiva}
\newcommand{\figroot}{../coinfosim-report-latex/figures}

% Linha de fonte, sempre no rodapé do conteúdo, em cinza discreto.
\newcommand{\sourcecite}[1]{%
  \par\vspace{0.05cm}%
  {\color{CoInfoGray}\fontsize{6.6}{7.8}\selectfont Fonte: #1\par}%
}

% Rótulo curto em caixa alta acima de um bloco de conteúdo.
\newcommand{\eyebrow}[1]{%
  {\bfseries\fontsize{8.3}{9.8}\selectfont\color{CoInfoBlue}\MakeUppercase{#1}}%
}

% Mensagem principal do slide: caixa creme, texto em negrito.
\newcommand{\takeaway}[1]{%
  \begin{beamercolorbox}[wd=\linewidth,sep=0.14cm,rounded=true]{takeaway}
    \centering\bfseries\fontsize{10.4}{12.3}\selectfont #1
  \end{beamercolorbox}%
}

% Cartão informativo: título curto em destaque e corpo explicativo completo.
\newcommand{\infocard}[3][\linewidth]{%
  \begin{beamercolorbox}[wd=#1,sep=.12cm,rounded=true]{block body}
    {\bfseries\fontsize{9.8}{11.6}\selectfont #2\par}
    \vspace{0.03cm}
    {\fontsize{8.7}{10.4}\selectfont #3\par}
  \end{beamercolorbox}%
}

% Pílula de destaque com borda suave.
\newcommand{\metricpill}[3]{%
  \begin{tikzpicture}[baseline=(n.base)]
    \node[rounded corners=2.6pt,draw=#1!55,fill=#1!7,inner xsep=5.5pt,inner ysep=3.4pt,
      text=CoInfoInk,font=\fontsize{8.6}{10}\selectfont] (n)
      {\textbf{#2}\; #3};
  \end{tikzpicture}%
}

% Cartão de resultados usado no material de apoio.
\newcommand{\resultcard}[4]{%
  \begin{beamercolorbox}[wd=\linewidth,sep=0.15cm,rounded=true]{block body}
    {\bfseries\color{#1}\fontsize{10}{12}\selectfont #2}\par\vspace{0.06cm}
    {\fontsize{9}{11}\selectfont #3}\par
    {\bfseries\fontsize{9}{11}\selectfont #4}
  \end{beamercolorbox}%
}

\newcommand{\speakernotes}[5]{%
  \note{%
    \fontsize{8.0}{9.5}\selectfont
    \textbf{Roteiro oral.} #1\par\smallskip
    \textbf{Transição.} #2\par\smallskip
    \textbf{Frase para lembrar.} #3\par\smallskip
    \textbf{Tempo planejado.} #4\par\smallskip
    \textbf{Cuidado conceitual.} #5
  }%
}

\newcommand{\backupbegin}{%
  \appendix
  \backupslidestrue
  \slidetime{}
}
